\hypertarget{introduction}{%
\subsection{Introduction}\label{introduction}}

To complete this activity, I began by going through a list of recently
published CVEs. This was done using NIST's National Vulnerability
Database, which provides a helpful search tool with filtering
capabilities making it easy to locate recently published
vulnerabilities. After some scrolling, I came across a few CVEs which
looked interesting but I wanted to use one where the vulnerability had
publicly accessible code that I could reference for the LLMs to actually
fix since that would give the most varied response (at least in theory).
Eventually I settled on CVE-2026-35363, which was published on April 22
2026 and impacted the \texttt{rm} implementation found in the uutils
coreutils library.

The uutils coreutils library is a reimplementation of the GNU coreutils
written in Rust, primarily being used to replace the C-based GNU
coreutils in Linux-based operating systems. This vulnerability was found
by Canonical, who audited the uutils coreutils implementations for
security flaws before they could be used to replace the GNU coreutils
currently used by Ubuntu. This audit led to 113 issues being identified,
of which this was one of them. Although many of these issues were
resolved before Ubuntu 26.04 LTS was released (and hence are currently
being used as of writing this), the \texttt{cp}, \texttt{mv}, and
\texttt{rm} commands are still using the GNU coreutils implementations
with plans to replace them with the uutils coreutils implementations in
Ubuntu 26.10.

The bug identified with CVE-2026-35363 is a lack of proper safeguard
mechanisms for the uutils coreutils implementation of \texttt{rm}.
Specificially, it is vulnerable to CWE-22 (improper limitation of a
pathname to a restricted directory) making it vulnerable to path
traversal attacks. The \texttt{rm} command should, normally, prevent the
current directory (or the parent directory) from being removed. However,
this implementation only prevent specific paths from being removed
(i.e.~\texttt{.} and \texttt{..}). It did not consider semantically
equivalent paths, like \texttt{./} or \texttt{.././}, which could be
used to delete all contents within the current directory. The reason
this wasn't discovered earlier is because the command did report that it
received an invalid input, but deleted everything in the specified
directory regardless.

The three LLMs chosen were ChatGPT (using GPT-5.2), Claude (using Sonnet
4.6 Adaptive), and Gemini (using Gemini 3 Flash). In each case, the free
tier was used (using the models available as of April 25 2026) to
minimise the difference in model quality between vendors (i.e.~it uses
the ``baseline'' models for each to give a true representation of their
capabilities) in addition to it being infeasible to pay for three
different subscriptions to access higher tier models. Although this will
give a baseline benchmark for LLMs, it should be noted that using higher
tier models could give more effective output.

The following prompt was given to all models with no alterations.
\textgreater{} Consider CVE-2026-35363. This is a vulnerability that
affects the \texttt{rm} utility of uutils coreutils, a Rust
reimplementation of the GNU coreutils. The following is a description of
the issue: \textgreater{}
\texttt{\textgreater{}\ A\ vulnerability\ in\ the\ rm\ utility\ of\ uutils\ coreutils\ allows\ the\ bypass\ of\ safeguard\ mechanisms\ intended\ to\ protect\ the\ current\ directory.\ While\ the\ utility\ correctly\ refuses\ to\ delete\ .\ or\ ..,\ it\ fails\ to\ recognize\ equivalent\ paths\ with\ trailing\ slashes,\ such\ as\ ./\ or\ .///.\ An\ accidental\ or\ malicious\ execution\ of\ rm\ -rf\ ./\ results\ in\ the\ silent\ recursive\ deletion\ of\ all\ contents\ within\ the\ current\ directory.\ The\ command\ further\ obscures\ the\ data\ loss\ by\ reporting\ a\ misleading\ \textquotesingle{}Invalid\ input\textquotesingle{}\ error,\ which\ may\ cause\ users\ to\ miss\ the\ critical\ window\ for\ data\ recovery.\ \textgreater{}}
\textgreater{} \textgreater{} This issue is identified by issue \#9749
on the GitHub page for this library, with the following description:
\textgreater{}
\texttt{\textgreater{}\ Component\ \textgreater{}\ rm\ -rf\ \textgreater{}\ \ \textgreater{}\ Description\ \textgreater{}\ An\ issue\ with\ rm\ -rf:\ when\ deleting\ with\ rm\ -rf\ .,\ consistent\ with\ GNU\ rm,\ it\ refuses\ to\ delete\ the\ current\ directory,\ reports\ an\ error\ directly,\ and\ does\ not\ delete\ any\ contents.\ \textgreater{}\ \ \textgreater{}\ rm\ -rf\ .\ \textgreater{}\ rm:\ refusing\ to\ remove\ \textquotesingle{}.\textquotesingle{}\ or\ \textquotesingle{}..\textquotesingle{}\ directory:\ skipping\ \textquotesingle{}.\textquotesingle{}\ \textgreater{}\ However,\ when\ using\ rm\ -rf\ ./\ or\ rm\ -rf\ .///,\ it\ reports\ an\ error\ as\ if\ the\ deletion\ failed\ (it\ should\ be\ consistent\ with\ rm\ -rf\ .,\ refusing\ directly\ without\ deleting\ anything),\ but\ in\ reality\ the\ contents\ of\ the\ current\ directory\ have\ already\ been\ deleted:\ \textgreater{}\ \ \textgreater{}\ rm:\ cannot\ remove\ \textquotesingle{}./\textquotesingle{}:\ Invalid\ input\ \textgreater{}\ This\ is\ because\ the\ clean\_trailing\_slashes\ function\ collapses\ .///\ into\ "./":\ \textgreater{}\ \ \textgreater{}\ fn\ clean\_trailing\_slashes(path:\ \&Path)\ -\textgreater{}\ \&Path\ \{\ \textgreater{}\ \ \ \ \ ...\ \textgreater{}\ \ \ \ \ \ \ \ \ //\ Checks\ if\ element\ at\ the\ end\ is\ a\ \textquotesingle{}/\textquotesingle{}\ \textgreater{}\ \ \ \ \ \ \ \ \ if\ path\_bytes{[}idx{]}\ ==\ dir\_separator\ \{\ \textgreater{}\ \ \ \ \ \ \ \ \ \ \ \ \ for\ i\ in\ (1..path\_bytes.len()).rev()\ \{\ \textgreater{}\ \ \ \ \ \ \ \ \ \ \ \ \ \ \ \ \ if\ path\_bytes{[}i\ -\ 1{]}\ !=\ dir\_separator\ \{\ \textgreater{}\ \ \ \ \ \ \ \ \ \ \ \ \ \ \ \ \ \ \ \ \ idx\ =\ i;\ \textgreater{}\ \ \ \ \ \ \ \ \ \ \ \ \ \ \ \ \ \ \ \ \ break;\ \textgreater{}\ \ \ \ \ \ \ \ \ \ \ \ \ \ \ \ \ \}\ \textgreater{}\ \ \ \ \ \ \ \ \ \ \ \ \ \}\ \textgreater{}\ \ \ \ \ \ \ \ \ \ \ \ \ \#{[}cfg(unix){]}\ \textgreater{}\ \ \ \ \ \ \ \ \ \ \ \ \ return\ Path::new(OsStr::from\_bytes(\&path\_bytes{[}0..=idx{]}));\ \textgreater{}\ \ \ \ \ \ \ \ \ \ \ \ \ ...\ \textgreater{}\ \ \ \ \ \ \ \ \ \}\ \textgreater{}\ \ \ \ \ \}\ \textgreater{}\ \ \ \ \ path\ \textgreater{}\ \}\ \textgreater{}\ However,\ the\ path\_is\_current\_or\_parent\_directory\ function\ only\ matches\ ".",\ "..",\ etc.\ It\ does\ not\ match\ equivalent\ paths\ like\ "./"\ or\ "../".\ \textgreater{}\ \ \textgreater{}\ ///\ Checks\ if\ the\ path\ is\ referring\ to\ current\ or\ parent\ directory,\ if\ it\ is\ referring\ to\ current\ or\ any\ parent\ directory\ in\ the\ file\ tree\ e.g\ \textquotesingle{}/../..\textquotesingle{},\ \textquotesingle{}../..\textquotesingle{}\ \textgreater{}\ fn\ path\_is\_current\_or\_parent\_directory(path:\ \&Path)\ -\textgreater{}\ bool\ \{\ \textgreater{}\ \ \ \ \ let\ path\_str\ =\ os\_str\_as\_bytes(path.as\_os\_str());\ \textgreater{}\ \ \ \ \ let\ dir\_separator\ =\ MAIN\_SEPARATOR\ as\ u8;\ \textgreater{}\ \ \ \ \ if\ let\ Ok(path\_bytes)\ =\ path\_str\ \{\ \textgreater{}\ \ \ \ \ \ \ \ \ return\ path\_bytes\ ==\ ({[}b\textquotesingle{}.\textquotesingle{}{]})\ \textgreater{}\ \ \ \ \ \ \ \ \ \ \ \ \ \textbar{}\textbar{}\ path\_bytes\ ==\ ({[}b\textquotesingle{}.\textquotesingle{},\ b\textquotesingle{}.\textquotesingle{}{]})\ \textgreater{}\ \ \ \ \ \ \ \ \ \ \ \ \ \textbar{}\textbar{}\ path\_bytes.ends\_with(\&{[}dir\_separator,\ b\textquotesingle{}.\textquotesingle{}{]})\ \textgreater{}\ \ \ \ \ \ \ \ \ \ \ \ \ \textbar{}\textbar{}\ path\_bytes.ends\_with(\&{[}dir\_separator,\ b\textquotesingle{}.\textquotesingle{},\ b\textquotesingle{}.\textquotesingle{}{]})\ \textgreater{}\ \ \ \ \ \ \ \ \ \ \ \ \ \textbar{}\textbar{}\ path\_bytes.ends\_with(\&{[}dir\_separator,\ b\textquotesingle{}.\textquotesingle{},\ dir\_separator{]})\ \textgreater{}\ \ \ \ \ \ \ \ \ \ \ \ \ \textbar{}\textbar{}\ path\_bytes.ends\_with(\&{[}dir\_separator,\ b\textquotesingle{}.\textquotesingle{},\ b\textquotesingle{}.\textquotesingle{},\ dir\_separator{]});\ \textgreater{}\ \ \ \ \ \}\ \textgreater{}\ \ \ \ \ false\ \textgreater{}\ \}\ \textgreater{}}
\textgreater{} \textgreater{} How would you fix this particular
vulnerability, and why would you make that change?

Note: the code generated by each model will not be copied here. The
chats have been provided for your own viewing, where the exact code
generated can be seen. This will also not be tested because the codebase
is too large to reasonably test the changes, and I don't know Rust well
enough to be able to run the command properly. Instead, the conclusions
each model comes to will be summarised and compared which should
hopefully be enough.

\hypertarget{chatgpt-gpt-5.2}{%
\subsubsection{\texorpdfstring{\href{https://chatgpt.com/share/69ec35c1-2c74-839e-8fa2-288cdefbc70d}{ChatGPT
(GPT-5.2)}}{ChatGPT (GPT-5.2)}}\label{chatgpt-gpt-5.2}}

ChatGPT identified that the cause of the bug was a lack of proper safety
checks. To fix this bug, it recommended a few different options: -
Normalising the path components before checking - Explicitly include
semantically equivalent paths in
\texttt{path\_is\_current\_or\_parent\_directory}

The normalising approach uses the \texttt{Path} module in Rust's
standard library, which has a \texttt{components} method which
automatically normalises paths. This would allow it to handle all
equivalent representations of the same directory, ensuring that this
function doesn't silently fail.

The explicit path inclusion approach involves adding paths like
\texttt{./} and \texttt{../} to the function, but it is noted that this
is a weaker approach that will keep missing edge cases (as these paths
can be arbitrarily long, such as \texttt{././{[}...{]}./}).

\hypertarget{claude-sonnet-4.6-adaptive}{%
\subsubsection{\texorpdfstring{\href{https://claude.ai/share/7e50c827-5c7d-4b5b-967f-08524e29c423}{Claude
(Sonnet 4.6
Adaptive)}}{Claude (Sonnet 4.6 Adaptive)}}\label{claude-sonnet-4.6-adaptive}}

Claude identified that the cause of the bug was
\texttt{clean\_trailing\_slashes} strips trailing slahes but leaves one
which is then used with
\texttt{path\_is\_current\_or\_parent\_directory}. It proposes a few
different options to fix this bug: - Strip all trailing slahes in
\texttt{clean\_trailing\_slashes} - Add the missing \texttt{./} and
\texttt{../} patterns to
\texttt{path\_is\_current\_or\_parent\_directory}

The stripping trailing slashes approach involves removing all trailing
slashes from the path string so that the path always ends in a non
\texttt{/} character. This proposes this because it will normalise the
provided paths, which then gets caught by
\texttt{path\_is\_current\_or\_parent\_directory}.

The adding missing patterns approach involves just adding 2 missing
paths - \texttt{./} and \texttt{../}. It suggests this approach because
the \texttt{path\_is\_current\_or\_parent\_directory} should recognise
any path that refers to the current or parent directory, even if the
input isn't normalised.

It recommends using both of these approaches together because each
approach misses something if it is implemented alone.

\hypertarget{gemini-gemini-3-flash}{%
\subsubsection{\texorpdfstring{\href{https://gemini.google.com/share/eb2352e5b683}{Gemini
(Gemini 3
Flash)}}{Gemini (Gemini 3 Flash)}}\label{gemini-gemini-3-flash}}

Gemini identified the cause of the bug was the safeguard being too
literal and the input path not being properly cleaned/normalised. The
fix it proposes is to normalise the path or update the
\texttt{path\_is\_current\_or\_parent\_directory} function to be aware
of the trailing \texttt{/} character. Gemini proposes this fix because
it fixes the logical flaw and makes behaviour consistent with GNU
coreutils.

\hypertarget{comparison}{%
\subsubsection{Comparison}\label{comparison}}

At the time of writing, no change has been made to this vulnerability
yet. However, pull request \#11005 has been created to address the
problem, which proposes the following solution:

\begin{Shaded}
\begin{Highlighting}[]
\KeywordTok{fn}\NormalTok{ path\_is\_current\_or\_parent\_directory(path}\OperatorTok{:} \OperatorTok{\&}\DataTypeTok{Path}\NormalTok{) }\OperatorTok{{-}\textgreater{}} \DataTypeTok{bool} \OperatorTok{\{}
    \KeywordTok{let} \ConstantTok{Ok}\NormalTok{(bytes) }\OperatorTok{=}\NormalTok{ os\_str\_as\_bytes(path}\OperatorTok{.}\NormalTok{as\_os\_str()) }\ControlFlowTok{else} \OperatorTok{\{}
        \ControlFlowTok{return} \ConstantTok{false}\OperatorTok{;}
    \OperatorTok{\};}

    \CommentTok{// Strip all trailing separators}
    \KeywordTok{let}\NormalTok{ trimmed }\OperatorTok{=} \ControlFlowTok{match}\NormalTok{ bytes}\OperatorTok{.}\NormalTok{iter()}\OperatorTok{.}\NormalTok{rposition(}\OperatorTok{|\&}\NormalTok{b}\OperatorTok{|} \OperatorTok{!}\NormalTok{is\_separator(b)) }\OperatorTok{\{}
        \ConstantTok{Some}\NormalTok{(i) }\OperatorTok{=\textgreater{}} \OperatorTok{\&}\NormalTok{bytes[}\OperatorTok{..=}\NormalTok{i]}\OperatorTok{,}
        \ConstantTok{None} \OperatorTok{=\textgreater{}} \ControlFlowTok{return} \ConstantTok{false}\OperatorTok{,} \CommentTok{// all separators or empty}
    \OperatorTok{\};}

    \CommentTok{// Extract the last component (after the last separator)}
    \KeywordTok{let}\NormalTok{ last }\OperatorTok{=} \ControlFlowTok{match}\NormalTok{ trimmed}\OperatorTok{.}\NormalTok{iter()}\OperatorTok{.}\NormalTok{rposition(}\OperatorTok{|\&}\NormalTok{b}\OperatorTok{|}\NormalTok{ is\_separator(b)) }\OperatorTok{\{}
        \ConstantTok{Some}\NormalTok{(i) }\OperatorTok{=\textgreater{}} \OperatorTok{\&}\NormalTok{trimmed[i }\OperatorTok{+} \DecValTok{1}\OperatorTok{..}\NormalTok{]}\OperatorTok{,}
        \ConstantTok{None} \OperatorTok{=\textgreater{}}\NormalTok{ trimmed}\OperatorTok{,}
    \OperatorTok{\};}

\NormalTok{    last }\OperatorTok{==} \StringTok{b"."} \OperatorTok{||}\NormalTok{ last }\OperatorTok{==} \StringTok{b".."}
\OperatorTok{\}}
\end{Highlighting}
\end{Shaded}

This is different to what each LLM proposed, and instead strips all the
trailing \texttt{/} characters and checks that the last component is
\texttt{.} or \texttt{..}. Although this is effectively just normalising
the path, the closest answer to this solution was provided by Claude.

The least surprising thing about the answers provided by each LLM is
that the primary solution was to normalise the path, which the proposed
solution also does. Additionally, each LLM makes reference to the
potential data loss vulnerability created from the output generated when
paths like \texttt{./} are passed. However, only ChatGPT and Claude
provide an alternative solution. This alternative solution is the same
between them, and involves explicitly adding missing paths to the final
check. However, ChatGPT specifically makes note that this is a weak
solution while Claude says that both solutions should be used together
rather than choosing between one or the other.

An interesting thing that deviates between the models is their
justification for their proposal. Both Claude and Gemini justify why
they chose to implement their solution in the way they chose, which use
reasonably sound logic. Gemini in particular gives several
justifications for its reasoning, though every reason is semi-related
(which can be forgiven). However, ChatGPT justifies why its
implementation works. Although this is just a semantic difference since
it ends up essentially being the same, its still an interesting shift
regardless.

\hypertarget{references}{%
\subsubsection{References}\label{references}}

National Institute of Standards and Technology. ``NVD Vulnerability
Search''. Accessed: Apr.~24, 2026. {[}Online{]}. Available:
https://nvd.nist.gov/vuln/search\#/nvd/home?vulnRevisionStatusList=Published\&sortOrder=3\&sortDirection=2\&resultType=records

National Institute of Standards and Technology. ``CVE-2026-35363
Detail''. Accessed: Apr.~25, 2026. {[}Online{]}. Available:
https://nvd.nist.gov/vuln/detail/CVE-2026-35363

R. Sharma. ``An update on rust-coreutils''. Canonical Ltd.~Accessed:
Apr.~25, 2026. {[}Online{]}. Available:
https://discourse.ubuntu.com/t/an-update-on-rust-coreutils/80773

The MITRE Corporation. ``CVE-2026-35363''. Accessed: Apr.~25, 2026.
{[}Online{]}. Available: https://www.cve.org/CVERecord?id=CVE-2026-35363

The MITRE Corporation. ``CWE-22: Improper Limitation of a Pathname to a
Restricted Directory (`Path Traversal')''. Accessed: Apr.~25, 2026.
{[}Online{]}. Available: https://cwe.mitre.org/data/definitions/22.html

The Rust Team. ``Module path''. Accessed: Apr.~25, 2026. {[}Online{]}.
Available: https://doc.rust-lang.org/std/path/index.html
